\color{unidarkgreen}{\subsection{Дистанционная защита}}
\color{black}

\begin{enumerate}
%%%%%%%%%%%%%%%%%%%%%%%%%%%%%%%%%%%%%%%%%%%%%%%%%% ОБЩИЕ СВЕДЕНИЯ %%%%%%%%%%%%%%%%%%%%%%%%%%%%%%%%%%%%%%%%%%%%%%%%%%
\item Общие сведения

\begin{enumerate}

\item
Дистанционная защита (ДЗ) предназначена для селективного отключения поврежденного элемента электрической сети при всех видах коротких замыканий (КЗ), реагируя на снижение измеряемого сопротивления в защищаемой зоне.

\item
Уставки ДЗ представлены в таблице \ref{app_set:dz}.

\item 
ДЗ селективно срабатывает:
\begin{itemize}
  \item при всех видах КЗ, в том числе при постановке линии под напряжение; 
  \item при переходе внешнего КЗ во внутреннее, в том числе при реверсе мощности; 
  \item при КЗ, сопровождающемся снижением измеряемого напряжения до нуля.
\end{itemize}

\item 
Обеспечивается несрабатывание защиты:
\begin{itemize} 
	\item при качаниях и асинхронном ходе; 
	\item при неисправностях цепей напряжения;
	\item при внешних КЗ с насыщением ТТ (при соблюдении заявленных производителем требований к ТТ); 
	\item при постановке линии под напряжение и включении в транзит в отсутствие КЗ;
	\item при реверсе мощности при внешних КЗ; 
	\item при бросках тока намагничивания без КЗ;
	\item при внешних КЗ, даже если измеряемое напряжение снижается до нуля.
\end{itemize}

\item Каждая из ступеней ДЗ от междуфазных КЗ включает в себя три реле сопротивления (РС), использующие для замера сопротивления комплексные значения линейных напряжений и соответствующих им линейных токов:

\begin{align}
\dot Z_{AB} = \frac{\dot U_A - \dot U_B}{\dot I_A - \dot I_B}\label{eq:groupZ1}, \\[10pt]
\dot Z_{BC} = \frac{\dot U_B - \dot U_C}{\dot I_B - \dot I_C}\label{eq:groupZ2}, \\[10pt]
\dot Z_{CA} = \frac{\dot U_C - \dot U_A}{\dot I_C - \dot I_A}\label{eq:groupZ3}, 
\end{align}
\begin{eqexpl}[20mm]
  \item{$ \dot U_{A}, \dot U_{B}, \dot U_{C} $} комплексные значения напряжений фаз А, В, С;
	\vspace{2pt}
  \item{$ \dot I_{A}, \dot I_{B}, \dot I_{C} $} комплексные значения токов фаз А, В, С.
\end{eqexpl}

\item 
Комплексные значения сопротивлений по контуру «фаза-земля» рассчитываются с учетом компенсации тока нулевой последовательности своей линии и параллельной линии:

\begin{align}
\dot Z_{A} = \frac{\dot U_A }{\dot I_A + \dot K \cdot 3 \dot I_0 + \dot K_{\textsl{пл}} \cdot 3 \dot I_{0\textsl{ пл}}}\label{eq:groupZ11}, \\[10pt]
\dot Z_{B} = \frac{\dot U_B }{\dot I_B + \dot K \cdot 3 \dot I_0 + \dot K_{\textsl{пл}} \cdot 3 \dot I_{0\textsl{ пл}}}\label{eq:groupZ12}, \\[10pt]
\dot Z_{C} = \frac{\dot U_C }{\dot I_C + \dot K \cdot 3 \dot I_0 + \dot K_{\textsl{пл}} \cdot 3 \dot I_{0\textsl{ пл}}}\label{eq:groupZ13}, 
\end{align}
\begin{eqexpl}[5mm]
  \item{$ 3 \dot I_{0} $} комплексное значение утроенного тока нулевой последовательности защищаемой линии (расчетное значение);
  \item{$ 3 \dot I_{0\textsl{ пл}} $}  комплексное значение утроенного тока нулевой последовательности параллельной линии (измеренное значение);
  \item{$ \dot K $} комплексный коэффициент компенсации тока нулевой последовательности защищаемой линии;
  \item{$ \dot K_{\textsl{пл}} $}  комплексный коэффициент компенсации тока нулевой последовательности параллельной линии.
\end{eqexpl}

\item 
В устройстве коэффициент $\dot K$ задается двумя уставками, \textbf{«KR 3I0 своей линии»} и \textbf{«KХ 3I0 своей линии»}. коэффициент $ \dot K_{\textsl{пл}}$  задается уставками \textbf{«KR// 3I0 пар. линии»} и \textbf{«KХ// 3I0 пар. линии»}. 

\item 
Расчет коэффициентов компенсации зависит от режима работы параллельной линии. Если параллельная линия в работе, отключена или отсутствует, то коэффициенты компенсации рассчитываются по следующим формулам:

\begin{align}
\dot K = \frac{1}{3} \cdot \frac{\dot Z_{0 \textsl{ уд}} - \dot Z_{1 \textsl{ уд}}}{\dot Z_{1 \textsl{ уд}}}\label{eq:groupK1}, \\[10pt]
\dot K_{\textsl{пл}} = \frac{1}{3} \cdot \frac{\dot Z_{M \textsl{ уд}}}{\dot Z_{1 \textsl{ уд}}}\label{eq:groupK2}, 
\end{align}

\begin{eqexpl}[10mm]
	\item{$ \dot Z_{0 \textsl{ уд}} $} комплексное удельное сопротивление нулевой последовательности;
	\item{$ \dot Z_{1 \textsl{ уд}} $} комплексное удельное сопротивление прямой последовательности;
	\item{$ \dot Z_{M \textsl{ уд}} $} комплексное удельное сопротивление взаимоиндукции.
\end{eqexpl}

\item 
Компенсация влияния тока параллельной линии блокируется, если величина тока нулевой последовательности параллельной линии $ 3 \dot I_{0\textsl{ пл}} $ превышает 135\% от величины тока нулевой последовательности защищаемой линии $ 3 \dot I_{0} $:

\begin{equation}
3 \dot I_{0\textsl{ пл}} >1,35 \cdot 3 \dot I_{0}
\end{equation}

\item
Характеристики срабатывания ненаправленного РС для междуфазных и фазных контуров представляют собой параллелограмм, верхняя сторона которого параллельна оси R и пересекает ось Х в точке с координатой Хср, а правая сторона -- имеет угол наклона $φ_1$, соответствующий углу наклона характеристики, относительно оси R и пересекает ее в точке с координатой Rср (Хср и Rср - уставки соответствующей ступени по реактивному и активному сопротивлениям).

\item 
Вид характеристики срабатывания ненаправленного РС представлен на рисунке \ref{pict_dz:char1}. 

\medskip

\begin{figure}[H]
\centering
\includegraphics[width=0.45\textwidth,height=0.45\textheight,keepaspectratio]{\fbpath/03.01 ЛВ/561.0021 ДЗ/img/_char1.pdf}
\captionof{figure}{ Характеристика срабатывания ненаправленного РС}\label{pict_dz:char1}
\end{figure}

\item 
Характеристики срабатывания ненаправленного РС для междуфазных и фазных контуров задаются независимо для каждой ступени уставками:
\begin{itemize} 
	\item \textbf{«Rср ДЗ-ФФ(ФЗ)-N»} – активное сопротивление срабатывания РС; 
	\item \textbf{«Хср ДЗ-ФФ(ФЗ)-N»} – реактивное сопротивление срабатывания РС;
	\item \textbf{«Ф1 ДЗ-ФФ(ФЗ)-N»} – угол наклона характеристики (рекомендуется задавать равным характеристическому углу линии);
\end{itemize}

\begin{eqexpl}[3mm]
  \item{$ N $} номер ступени ДЗ.
\end{eqexpl} 

\item 
Для предотвращения неселективного срабатывания РС-ФЗ-1 и РС-ФФ-1 при внешних повреждениях через переходное сопротивление предусмотрена возможность отсечения области характеристики срабатывания, задаваемой углом $φ_4$. Область отсекается отрезком, исходящей под углом $φ_4$, из точки на верхней части характеристики, соответствующей углу линии электропередачи $φ_1$. Угол $φ_4$ задается уставкой \textbf{«Ф4 ДЗ-ФФ(ФЗ)-1»}. 

\item 
Параметры реле сопротивления приведены в таблице \ref{desc_dz:tbl1}

\small
\setlength{\extrarowheight}{0.05cm} % высота ячеек в таблице
\begin{longtable}
{|>{\centering\arraybackslash}m{15.2cm}|>{\centering\arraybackslash}m{1.7cm}|}
\caption{Параметры РС\hfill\vspace{-0.5\baselineskip}}\label{desc_dz:tbl1}\\ 
\hhline{|-|-|}
\rowcolor{gray!30}
Наименование параметра & Значение \\ 
\hhline{|-|-|}
\endfirsthead

\caption*{\hspace{3pt}\emph{Продолжение таблицы \ref{desc_dz:tbl1}\hfill\vspace{-0.5\baselineskip}}} \\ % сделано по ГОСТ 2.105 п.6.8.7
\hhline{|-|-|}
\rowcolor{gray!30}
Наименование параметра & Значение \\ 
\hhline{|-|-|}
\endhead
\hhline{|-|-|}

\endfoot
\endlastfoot

\raggedright Минимальное напряжение, при котором обеспечивается средняя основная погрешность срабатывания ИО РС по величине сопротивления срабатывания, В вторичных, не более 						   					& \centering\arraybackslash 1 \\ \hline
\raggedright Ток десятипроцентной точности работы РС при работе на угле линии электропередачи, \% от $I_{\textsl{ном}}$, не более 																	   					& \centering\arraybackslash 10 \\ \hline
\raggedright Средняя основная погрешность РС по величине сопротивления срабатывания $R_{\textsl{уст}}$ и $X_{\textsl{уст}}$, \% от уставки, не более 												   					& \centering\arraybackslash $\pm$5 \\ \hline
\raggedright Собственное время срабатывания РС при угле, равном углу линии, трехкратном токе точной работы и скачкообразном уменьшении напряжения в два раза по отношению к номинальному, мс, не более 					& \centering\arraybackslash 25 \\ \hline
\raggedright Собственное время возврата РС при угле, равном углу линии, трехкратном токе точной работы и скачкообразном увеличении напряжения от 0,1 до 1,0 номинального, мс, не более 				   					& \centering\arraybackslash 50 \\ \hline
\raggedright Дополнительная относительная погрешность при изменении частоты входного сигнала на ±5 Гц от номинальной не более, \% по отношению к значению при номинальной частоте				  	   					& \centering\arraybackslash $\pm$5 \\ \hline
\raggedright Дополнительная относительная погрешность при изменении температуры окружающей среды в диапазоне, указанном в таблице \ref{app_set:dz}, \% от значения, определенного при температуре (20±5) ºС, не более    & \centering\arraybackslash $\pm$5 \\ \hline
\raggedright Коэффициент возврата, не более  																																											& \centering\arraybackslash 1,1 \\ \hline
\raggedright Время срабатывания измерительного органа реле сопротивления при работе на угле линии электропередач, токе не менее 0,3$\cdot I_{\textsl{ном}}$ и скачкообразном уменьшении напряжения на входе реле от напряжения 100 В, соответствующего сопротивлению на входе реле не менее 1,2$\cdot(X_{\textsl{уст}} /sinφ_1 )$ до напряжения соответствующего 0,6$\cdot(X_{\textsl{уст}} /sinφ_1 )$, мс, не более 													& \centering\arraybackslash 25 \\ \hline
\raggedright Время возврата измерительного органа реле сопротивления при работе на угле линии электропередач, токе не менее 0,3$\cdot I_{\textsl{ном}}$ и скачкообразном увеличении напряжения на входе реле от напряжения соответствующего сопротивлению на входе реле 0,1$\cdot(X_{\textsl{уст}} /sinφ_1 )$ до напряжения соответствующего 1,2$\cdot(X_{\textsl{уст}} /sinφ_1 )$ (но не более 100 В для междуфазного канала и 57 В – для фазного), мс, не более		& \centering\arraybackslash 50 \\ \hline
\raggedright Время срабатывания ДЗ при переходе внешнего КЗ во внутреннее в условиях наличия насыщения ТТ мс, не более.																									& \centering\arraybackslash 60 \\ \hline

\end{longtable}
\setlength{\extrarowheight}{0.0cm} % возврат параметра высота ячеек в таблице
\normalsize

%%%%%%%%%%%%%%%%%%%%%%%%%%%%%%%%%%%%%%%%%%%%%%%%%% ОН %%%%%%%%%%%%%%%%%%%%%%%%%%%%%%%%%%%%%%%%%%%%%%%%%%

\item 
Выполнение функций направленного РС обеспечивается совместным использованием ненаправленного РС и органов направленности (ОН). Орган направленности представляет собой два отрезка, исходящие из начала координат и расположенные во втором и четвёртом квадрантах. Вид характеристики определяется углами $φ_2$ и $φ_3$ наклона этих отрезков, отсчитываемых относительно оси $R$. Углы наклона характеристики задаются уставками \textbf{«Ф2 ОН-ФФ(ФЗ)»} и\textbf{«Ф3 ОН-ФФ(ФЗ)»}. Уставки органов направленности являются общими для всех ступеней, задаются раздельно для контуров «фаза-фаза» и «фаза-земля». Углы наклона в устройстве задаются для прямонаправленного ОН, а характеристика обратнонаправленного ОН симметрична характеристике прямонаправленного ОН с поворотом на 180 градусов.

\item
Характеристики срабатывания органов направленности представлены на рисунке \ref{pict_dz:_on}. 

\medskip

\begin{figure}[H]
\centering
\includegraphics[width=0.65\textwidth,height=0.65\textheight,keepaspectratio]{\fbpath/03.01 ЛВ/561.0021 ДЗ/img/_on.pdf}
\captionof{figure}{Характеристики срабатывания органов направленности}\label{pict_dz:_on}
\end{figure}

\medskip

\item
Направленность для каждой ступени ДЗ задается с помощью программных переключателей \textbf{«Направл. ДЗ-ФФ(ФЗ)-N»}, имеющих три положения:
\begin{itemize}
	\item \textit{Ненапр.} – ступень выполнена ненаправленной; 
	\item \textit{Прямо} – ступень выполнена прямонаправленной; 
	\item \textit{Обратно} – ступень выполнена обратнонаправленной. 
\end{itemize}

\item
При близких трехфазных КЗ, когда все напряжения на входе РС близки к нулю, для определения направленности в течение времени не менее 40 мс используются напряжения предаварийного режима (работа по «памяти»). При работе РС «по памяти» при трехфазных КЗ в месте установки защиты обеспечивается длительность сигнала срабатывания на выходе РС не менее 40 мс в диапазоне токов от 0,2$\cdot I_{\textsl{ном}}$ до 30$\cdot I_{\textsl{ном}}$. При КЗ «за спиной» с токами до 30$\cdot I_{\textsl{ном}}$ обеспечивается отсутствие ложных срабатываний РС, в том числе при снижении измеряемого напряжения до нуля. 

\item
Функциональная схема логики органа направления мощности приведена на рисунке \ref{pict_dz:on}.

\medskip

\begin{figure}[H]
\centering
\includegraphics[width=0.55\textwidth,height=0.55\textheight,keepaspectratio]{\fbpath/03.01 ЛВ/561.0021 ДЗ/img/ON.pdf}
\captionof{figure}{Функциональная схема логики отстройки от нагрузочного режима для контуров ФФ и ФЗ}\label{pict_dz:on}
\end{figure}

\medskip

%%%%%%%%%%%%%%%%%%%%%%%%%%%%%%%%%%%%%%%%%%%%%%%%%% ВЫРЕЗ %%%%%%%%%%%%%%%%%%%%%%%%%%%%%%%%%%%%%%%%%%%%%%%%%%

\item
Для отстройки от нагрузочного режима возможен вырез в характеристике срабатывания РС. Измерительный орган отстройки от нагрузки используется для всех ступеней. Уставки задаются раздельно для фазных и междуфазных контуров. Характеристика срабатывания ИО приведена на рисунке \ref{pict_dz:r}. Эта характеристика срабатывания определяется уставками: 
\begin{itemize}
  \item \textbf{«Rнг ДЗ-ФФ(ФЗ)»} - сопротивление выреза нагрузочного режима; 
  \item \textbf{«Фнг ДЗ-ФФ(ФЗ)»} - угол выреза нагрузочного режима.   
\end{itemize}

\item
Исключаемая область в характеристике срабатывания симметрична относительно осей координат. Если нет необходимости отстройки от нагрузки, то \textbf{«Rнг ДЗ-ФФ(ФЗ)»} выбирается больше уставки по $R$ самой чувствительной ступени ДЗ. 

\medskip

\begin{figure}[H]
\centering
\includegraphics[width=0.45\textwidth,height=0.45\textheight,keepaspectratio]{\fbpath/03.01 ЛВ/561.0021 ДЗ/img/_r.pdf}
\captionof{figure}{Характеристика срабатывания измерительного органа отстройки от нагрузки}\label{pict_dz:r}
\end{figure}

\medskip

\item
Функциональная схема логики отстройки от нагрузочного режима приведена на рисунке \ref{pict_dz:Zng}. 

\item
Уставки органов направленности и измерительных органов отстройки от нагрузки представлены в таблице \ref{app_set:dz}.

\begin{figure}[H]
\centering
\includegraphics[width=0.6\textwidth,height=0.6\textheight,keepaspectratio]{\fbpath/03.01 ЛВ/561.0021 ДЗ/img/Zng.pdf}
\captionof{figure}{Функциональная схема логики отстройки от нагрузочного режима для контуров ФФ и ФЗ}\label{pict_dz:Zng}
\end{figure}

\item
Суммарная характеристика направленного РС, используемого в логике работы ступеней ДЗ приведена на рисунке \ref{pict_dz:rs}. 

\begin{figure}[H]
\centering
\includegraphics[width=0.65\textwidth,height=0.65\textheight,keepaspectratio]{\fbpath/03.01 ЛВ/561.0021 ДЗ/img/_rs.pdf}
\captionof{figure}{Характеристики срабатывания направленного РС}\label{pict_dz:rs}
\end{figure}

\item
Дополнительно для каждой ступени ДЗ предусмотрены ненаправленные РС с охватом начала координат, уставки по $R$, $X$ и $φ_1$ которых совпадают с уставками ненаправленного РС соответствующей ступени. Характеристика срабатывания (см. рисунок \ref{pict_dz:nachalo}) ненаправленного РС с охватом начала координат имеет форму параллелограмма, смещенного вдоль направления металлического КЗ на линии в третий и четвертый квадранты на величину до 0,1$\cdot X_{\textsl{ср}}$. 

\begin{figure}[H]
\centering
\includegraphics[width=0.35\textwidth,height=0.35\textheight,keepaspectratio]{\fbpath/03.01 ЛВ/561.0021 ДЗ/img/_nachalo.pdf}
\captionof{figure}{Характеристика срабатывания ненаправленного РС с охватом начала координат}\label{pict_dz:nachalo}
\end{figure}

%%%%%%%%%%%%%%%%%%%%%%%%%%%%%%%%%%%% ООВП  %%%%%%%%%%%%%%%%%%%%%%%%%%%%%%%%%%

\item
Для обнаружения однофазных КЗ и пуска ступеней ДЗ-ФЗ, контролирующих контуры «фаза-земля», предусмотрен орган определения вида повреждения (ООВП), основанный на совместном применении измерительного органа тока нулевой последовательности (ИО 3I0 ООВП) с торможением от одного из фазных токов (средний из $I_A , I_B , I_C$) и измерительного органа напряжения нулевой последовательности (ИО 3U0 ООВП). Торможение осуществляется от тока той фазы, значение в которой является средним между максимальным и минимальным значениями токов в остальных двух фазах. В таком случае при междуфазном КЗ на землю торможение максимальном, а при однофазном КЗ –- минимально. 

\item
Характеристика срабатывания ИО 3I0 ООВП с торможением представлена на рисунке \ref{pict_dz:torm}.

\medskip
\begin{figure}[H]
\centering
\includegraphics[width=0.6\textwidth,height=0.6\textheight,keepaspectratio]{\fbpath/03.01 ЛВ/561.0021 ДЗ/img/_torm.pdf}
\captionof{figure}{Характеристика срабатывания ИО 3I0 ООВП}\label{pict_dz:torm}
\end{figure}
\medskip

\item
Порог срабатывания горизонтального участка задается уставкой \textbf{«3I0ср ООВП»}. Наклонный участок характеристики проходит через точку, равной $1,25 \cdot I_{\textsl{ном}}$ , под углом к горизонтальной оси, соответствующим коэффициенту торможения $k_T$. Коэффициент торможения задается уставкой \textbf{«Кт ООВП»}. При превышении тормозного тока уставки \textbf{«Iср блок. торм.»} ООВП блокируется. Алгоритм работы логики ООВП приведен на рисунке \ref{pict_dz:OOVP}.

\medskip
\begin{figure}[H]
\centering
\includegraphics[width=0.8\textwidth,height=0.8\textheight,keepaspectratio]{\fbpath/03.01 ЛВ/561.0021 ДЗ/img/OOVP.pdf}
\captionof{figure}{Функциональная схема логики ООВП}\label{pict_dz:OOVP}
\end{figure}
\medskip

\item
Ввод в работу ступеней ДЗ, работающих по контуру «фаза-земля» может осуществляться с дополнительным контролем напряжения нулевой последовательности за счет ввода программного ключа \textbf{«Контр. 3U0»}. Напряжение срабатывания ИО 3U0 ООВП задается уставкой \textbf{«3U0ср ООВП»}. ИО 3U0 ООВП может работать по напряжению $3U_{\textsl{0 звезда}}$, рассчитанному на основе фазных напряжений, или по напряжению $3U_{\textsl{0 треуг.}}$  -- выбирается программным ключом \textbf{«Выбор 3U0»}, имеющим два соответствующих положения: \textsl{«Расчетн.»} и \textsl{«Измерен.»}.

\item
Переменная $3U_{\textsl{0 треуг.}}$ зависит от положения программного переключателя \textbf{«Схема подключения цепей разомкнутого треугольника»}. Данный переключатель позволяет учесть различные схемы соединения дополнительной обмотки ТН. Если переключатель выставлен в положение «$U_{\textsl{НК}}$», то переменной $3U_{\textsl{0 треуг.}}$ присваивается значение напряжения, измеренного по аналоговому входу «$U_{\textsl{ИФ}}(U_{\textsl{ИК}})$». Если программный переключатель {\textbf{«Схема подключения цепей треугольника»}} выставлен в положение «$U_{\textsl{НИ}}/U_{\textsl{ИК}}$», «$U_{\textsl{НФ}}/U_{\textsl{ФК}}$» или «$U_{\textsl{НИ}}/U_{\textsl{ИФ}}/U_{\textsl{ФК}}$», то переменная $3U_{\textsl{0 треуг.}}$ рассчитывается в соответствии с формулами, приведенными в приложении \ref{app:bnn}. Программный переключатель \textbf{«Выбор 3U0»} аналогично влияет на работу ОНМНП. 

\item
Устройство содержит шесть ступеней ДЗ от междуфазных замыканий с контролем трех контуров «фаза-фаза», а также три ступени с контролем контуров «фаза-земля» для защиты от однофазных КЗ. 

\end{enumerate}

%%%%%%%%%%%%%%%%%%%%%%%%%%%%%%%%%%%%%%%%%%%%%%%%%% ПЕРВАЯ СТУПЕНЬ %%%%%%%%%%%%%%%%%%%%%%%%%%%%%%%%%%%%%%%%%%%%%%%%%%
\item\label{sec:sec2} Первая ступень ДЗ от междуфазных КЗ (ДЗ-ФФ-1)
\begin{enumerate}

\item
Первая ступень ДЗ от междуфазных КЗ (ДЗ-ФФ-1) вводится в работу программным ключом \textbf{«ДЗ-ФФ-1»}. Ступень имеет две регулируемые независимые выдержки времени действия на отключение: \textbf{«Тср ДЗ-ФФ-1.1»} и \textbf{«Тср ДЗ-ФФ-1.2»}. Действие с большей выдержкой времени вводится отдельным программным ключом \textbf{«ДЗ-ФФ-1.2»}. Контроль от блокировки при качаниях (БК) для каждой выдержки задается отдельными программными переключателями \textbf{«БК ДЗ-ФФ-1.1»} и \textbf{«БК ДЗ-ФФ-1.2»}.

\item
Для обеспечения надежного действия ступени при близких КЗ рекомендуется вводить подхват сработанного состояния РС-ФФ-1 от РС с охватом начала координат. Подхват вводится программным ключом \textbf{«Подхват РС-ФФ-1»}. Возврат схемы подхвата в исходное состояние происходит только после возврата направленного РС-ФФ-1 с охватом начала координат.

\item
Если междуфазное КЗ произошло в зоне действия ДЗ-ФФ-1 и при этом в течение времени ввода БК сработало направленное РС-ФФ-2 с охватом начала координат, то разрешающий сигнал от БК подхватывается и удерживается независимо от истечения времени ввода. Возврат разрешающего сигнала происходит только после возврата направленного РС-ФФ-2. 

\item
\textls[-5]{Функциональная схема логики первой ступени ДЗ от междуфазных КЗ приведена на рисунке \ref{pict_dz:ff1}.}

\medskip

\begin{figure}[H]
\centering
\includegraphics[width=0.85\textwidth,height=0.85\textheight,keepaspectratio]{\fbpath/03.01 ЛВ/561.0021 ДЗ/img/DZ_FF_1.pdf}
\captionof{figure}{Функциональная схема логики ДЗ-ФФ-1}\label{pict_dz:ff1}
\end{figure}

\medskip

\item
Уставки первой ступени ДЗ от междуфазных КЗ (ДЗ-ФФ-1) представлены в таблице \ref{app_set:dz}.

\end{enumerate}
%%%%%%%%%%%%%%%%%%%%%%%%%%%%%%%%%%%%%%%%%%%%%%%%%% КОНЕЦ ПЕРВАЯ СТУПЕНЬ %%%%%%%%%%%%%%%%%%%%%%%%%%%%%%%%%%%%%%%%%%%%%%%%%%


%%%%%%%%%%%%%%%%%%%%%%%%%%%%%%%%%%%%%%%%%%%%%%%%%% ВТОРАЯ СТУПЕНЬ %%%%%%%%%%%%%%%%%%%%%%%%%%%%%%%%%%%%%%%%%%%%%%%%%%
\item\label{sec:sec2} Вторая ступень ДЗ от междуфазных КЗ (ДЗ-ФФ-2)
\begin{enumerate}

\item
Вторая ступень ДЗ от междуфазных КЗ (ДЗ-ФФ-2) также выполнена с двумя выдержками времени и с раздельным контролем от БК. Вся ступень вводится программным ключом \textbf{«ДЗ-ФФ-2»}. Действие с меньшей выдержкой времени можно ввести в работу программным ключом \textbf{«ДЗ-ФФ-2.1»}. Выдержки времени действия на отключение регулируются уставками \textbf{«Тср ДЗ-ФФ-2.1»} и \textbf{«Тср ДЗ-ФФ-2.2»}. Контроль от БК для каждой выдержки задается отдельными программными переключателями \textbf{«БК ДЗ-ФФ-2.1»} и \textbf{«БК ДЗ-ФФ-2.2»}.

\item
Если КЗ произошло в зоне действия ступени и при этом в течение времени ввода БК сработало направленное РС, то разрешающий сигнал от БК подхватывается и удерживается независимо от истечения времени ввода. Возврат разрешающего сигнала происходит после возврата направленного РС.

\item
Функциональная схема логики второй ступени ДЗ от междуфазных КЗ приведена на рисунке \ref{pict_dz:ff2}.

\medskip

\begin{figure}[H]
\centering
\includegraphics[width=0.85\textwidth,height=0.85\textheight,keepaspectratio]{\fbpath/03.01 ЛВ/561.0021 ДЗ/img/DZ_FF_2.pdf}
\captionof{figure}{Функциональная схема логики ДЗ-ФФ-2}\label{pict_dz:ff2}
\end{figure}

\medskip

\item
Уставки второй ступени ДЗ от междуфазных КЗ (ДЗ-ФФ-2) представлены в таблице \ref{app_set:dz}.

\end{enumerate}

%%%%%%%%%%%%%%%%%%%%%%%%%%%%%%%%%%%%%%%%%%%%%%%%%% КОНЕЦ ВТОРАЯ СТУПЕНЬ %%%%%%%%%%%%%%%%%%%%%%%%%%%%%%%%%%%%%%%%%%%%%%%%%%


%%%%%%%%%%%%%%%%%%%%%%%%%%%%%%%%%%%%%%%%%%%%%%%%%% 3,4,5,6 СТУПЕНИ %%%%%%%%%%%%%%%%%%%%%%%%%%%%%%%%%%%%%%%%%%%%%%%%%%
\item\label{sec:sec21} Остальные ступени ДЗ от междуфазных КЗ (ДЗ-ФФ-3, ДЗ-ФФ-4, ДЗ-ФФ-5, ДЗ-ФФ-6)
\begin{enumerate}

\item
Логика работы остальных ступеней ДЗ от междуфазных КЗ выполнена одинаковой. Ввод в работу ступеней осуществляется программными ключами \textbf{«ДЗ-ФФ-3(4,5,6)»}. Действие ступеней осуществляется с регулируемыми независимыми выдержками времени \textbf{«Тср ДЗ-ФФ-3(4,5,6)»}. Контроль от БК задается отдельными программными переключателями \textbf{«БК ДЗ-ФФ-3(4,5,6)»}. 

\item
Разрешающий сигнал от БК для ступеней подхватывается от срабатывания РС соответствующей ступени и удерживается независимо от истечения времени ввода БК. Возврат разрешающего сигнала происходит только после возврата РС соответствующей ступени. На рисунке \ref{pict_dz:ff3} приведена функциональная схема логики ДЗ-ФФ-3, выполнение функциональных схем логики ступеней ДЗ ДЗ-ФФ-4, ДЗ-ФФ-5, ДЗ-ФФ-6 аналогичное.

\medskip

\begin{figure}[H]
\centering
\includegraphics[width=0.85\textwidth,height=0.85\textheight,keepaspectratio]{\fbpath/03.01 ЛВ/561.0021 ДЗ/img/DZ_FF_3.pdf}
\captionof{figure}{Функциональная схема логики ДЗ-ФФ-3}\label{pict_dz:ff3}
\end{figure}

\medskip

\item
Уставки ступеней ДЗ от междуфазных КЗ (ДЗ-ФФ-3, ДЗ-ФФ-4, ДЗ-ФФ-5, ДЗ-ФФ-6) представлены в таблице \ref{app_set:dz}.


\end{enumerate}
%%%%%%%%%%%%%%%%%%%%%%%%%%%%%%%%%%%%%%%%%%%%%%%%%% КОНЕЦ 3,4,5,6 СТУПЕНИ %%%%%%%%%%%%%%%%%%%%%%%%%%%%%%%%%%%%%%%%%%%%%%%%%%


%%%%%%%%%%%%%%%%%%%%%%%%%%%%%%%%%%%%%%%%%%%%%%%%%% ПЕРВАЯ СТУПЕНЬ ФАЗА-ЗЕМЛЯ %%%%%%%%%%%%%%%%%%%%%%%%%%%%%%%%%%%%%%%%%%%%%%%%%%
\item\label{sec:sec1fz} Первая ступень ДЗ от однофазных КЗ (ДЗ-ФЗ-1)
\begin{enumerate}

\item
Первая ступень ДЗ от однофазных КЗ (ДЗ-ФЗ-1) по аналогии с ДЗ-ФФ-1 выполнена с двумя выдержками времени и с раздельным контролем от БК. Вся ступень вводится программным ключом \textbf{«ДЗ-ФЗ-1»}. Действие с большей выдержкой времени можно ввести в работу программным ключом \textbf{«ДЗ-ФЗ-1.2»}. Выдержки времени действия на отключение регулируются уставками \textbf{«Тср ДЗ-ФЗ-1.1»} и \textbf{«Тср ДЗ-ФЗ-1.2»}. Контроль от БК для каждой выдержки задается отдельными программными переключателями \textbf{«БК ДЗ-ФЗ-1.1»} и \textbf{«БК ДЗ-ФЗ-1.2»}.

\item
Для обеспечения надежного действия ступени при близких КЗ рекомендуется вводить подхват сработанного состояния РС-ФЗ-1 от РС с охватом начала координат. Подхват вводится программным ключом \textbf{«Подхват РС-ФЗ-1»}. Возврат схемы подхвата в исходное состояние происходит только после возврата направленного РС-ФЗ-1 с охватом начала координат.

\item
Если КЗ произошло в зоне действия ступени и при этом в течение времени ввода БК сработало направленное РС, то разрешающий сигнал от БК подхватывается и удерживается независимо от истечения времени ввода. Возврат разрешающего сигнала происходит после возврата направленного РС.

\item
Функциональная схема логики первой ступени ДЗ от однофазных КЗ приведена на рисунке \ref{pict_dz:DZ_FZ_1}.

\medskip

\begin{figure}[H]
\centering
\includegraphics[width=0.85\textwidth,height=0.85\textheight,keepaspectratio]{\fbpath/03.01 ЛВ/561.0021 ДЗ/img/DZ_FZ_1.pdf}
\captionof{figure}{Функциональная схема логики ДЗ-ФЗ-1}\label{pict_dz:DZ_FZ_1}
\end{figure}

\medskip

\item
Уставки первой ступени ДЗ от однофазных КЗ (ДЗ-ФЗ-1) представлены в таблице \ref{app_set:dz}.

\end{enumerate}
%%%%%%%%%%%%%%%%%%%%%%%%%%%%%%%%%%%%%%%%%%%%%%%%%% КОНЕЦ ПЕРВАЯ СТУПЕНЬ ФАЗА-ЗЕМЛЯ %%%%%%%%%%%%%%%%%%%%%%%%%%%%%%%%%%%%%%%%%%%%%%%%%%


%%%%%%%%%%%%%%%%%%%%%%%%%%%%%%%%%%%%%%%%%%%%%%%%%% 4,5 СТУПЕНИ ФАЗА-ЗЕМЛЯ %%%%%%%%%%%%%%%%%%%%%%%%%%%%%%%%%%%%%%%%%%%%%%%%%%
\item\label{sec:secfz} Остальные ступени ДЗ от однофазных КЗ (ДЗ-ФЗ-4, ДЗ-ФЗ-5)
\begin{enumerate}

\item
Логика работы остальных ступеней ДЗ от однофазных КЗ выполнена одинаковой. Ввод в работу ступеней осуществляется программными ключами \textbf{«ДЗ-ФЗ-4(5)»}. Действие ступеней осуществляется с регулируемыми независимыми выдержками времени \textbf{«Тср ДЗ-ФЗ-4(5)»}. Контроль от БК задается отдельными программными переключателями \textbf{«БК ДЗ-ФЗ-4(5)»}.

\item
Разрешающий сигнал от БК для ступеней подхватывается от срабатывания РС соответствующей ступени и удерживается независимо от истечения времени ввода БК. Возврат разрешающего сигнала происходит только после возврата РС соответствующей ступени. На рисунке \ref{pict_dz:DZ_FZ_2} приведена функциональная схема логики ДЗ-ФЗ-4, выполнение функциональной схемы логики ступени ДЗ-ФЗ-5 аналогичное.

\medskip

\begin{figure}[H]
\centering
\includegraphics[width=0.9\textwidth,height=0.9\textheight,keepaspectratio]{\fbpath/03.01 ЛВ/561.0021 ДЗ/img/DZ_FZ_2.pdf}
\captionof{figure}{Функциональная схема логики ДЗ-ФЗ-4}\label{pict_dz:DZ_FZ_2}
\end{figure}

\medskip

\item
Уставки ступеней ДЗ от однофазных КЗ (ДЗ-ФЗ-4, ДЗ-ФЗ-5) представлены в таблице \ref{app_set:dz}.


\end{enumerate}
%%%%%%%%%%%%%%%%%%%%%%%%%%%%%%%%%%%%%%%%%%%%%%%%%% КОНЕЦ 4,5 СТУПЕНИ ФАЗА-ЗЕМЛЯ %%%%%%%%%%%%%%%%%%%%%%%%%%%%%%%%%%%%%%%%%%%%%%%%%%


%%%%%%%%%%%%%%%%%%%%%%%%%%%%%%%%%%%%%%%%%%%%%%%%%% БНН %%%%%%%%%%%%%%%%%%%%%%%%%%%%%%%%%%%%%%%%%%%%%%%%%%
\item\label{sec:secbnn} Взаимодействие с блокировкой при неисправностях цепей напряжения (БНН)
\begin{enumerate}

\item
Для исключения ложного действия при возникновении неисправности в цепях напряжения ступени ДЗ могут быть заблокированы от БНН. Блокировка от БНН вводится отдельно для каждой ступени ДЗ программными ключами \textbf{«БНН ДЗ-ФФ(ФЗ)-N»}.

\end{enumerate}
%%%%%%%%%%%%%%%%%%%%%%%%%%%%%%%%%%%%%%%%%%%%%%%%%% КОНЕЦ БНН %%%%%%%%%%%%%%%%%%%%%%%%%%%%%%%%%%%%%%%%%%%%%%%%%%


%%%%%%%%%%%%%%%%%%%%%%%%%%%%%%%%%%%%%%%%%%%%%%%%%% БК %%%%%%%%%%%%%%%%%%%%%%%%%%%%%%%%%%%%%%%%%%%%%%%%%%
\item\label{sec:secbk} Взаимодействие с блокировкой при качаниях (БК)
\begin{enumerate}

\item
Дистанционная защита контролируется блокировкой при качаниях, предотвращающей ложное срабатывание ДЗ при качаниях и асинхронном ходе. В устройстве реализовано два способа блокировки:
\begin{itemize}
  \item по приращению токов прямой и обратной последовательности;
  \item по скорости изменения величины сопротивления dZ/dt.
\end{itemize}

\item
Блокировка при качаниях, реагирующая на приращение во времени векторов тока прямой и обратной последовательностей, имеет два измерительных органа: чувствительный и грубый.

\item
Контроль ступеней ДЗ от блокировки при качаниях вводится программными переключателями \textbf{«БК ДЗ-ФФ(ФЗ)-N»}, имеющими четыре положения:
\begin{itemize}
	\item \textit{Откл.} -- ступень работает без контроля от БК;
	\item \textit{БК-I-б} -- ступень пускается от сигнала ввода быстродействующих ступеней;
	\item \textit{БК-I-м} -- ступень пускается от сигнала ввода медленнодействующих ступеней;
	\item \textit{БК-Z} -- ступень пускается от сигнала БК по скорости изменения величины сопротивления.
\end{itemize}

\end{enumerate}
%%%%%%%%%%%%%%%%%%%%%%%%%%%%%%%%%%%%%%%%%%%%%%%%%% КОНЕЦ БК %%%%%%%%%%%%%%%%%%%%%%%%%%%%%%%%%%%%%%%%%%%%%%%%%%


%%%%%%%%%%%%%%%%%%%%%%%%%%%%%%%%%%%%%%%%%%%%%%%%%% АУ ДЗ %%%%%%%%%%%%%%%%%%%%%%%%%%%%%%%%%%%%%%%%%%%%%%%%%%
\item\label{sec:secau} Автоматическое ускорение дистанционной защиты (АУ ДЗ)
\begin{enumerate}

\item
В составе ДЗ предусмотрена логика формирования пуска автоматического ускорения ДЗ (АУ ДЗ), выходной сигнал которого действует в логику автоматического ускорения. 
\item
Формирование сигнала пуска АУ ДЗ осуществляется от РС с охватом начала координат той ступени ДЗ-ФФ и ДЗ-ФЗ, которая выбрана программными переключателями \textbf{«Ступень АУ ДЗ-ФЗ»} и \textbf{«Ступень АУ ДЗ-ФФ»}. 

\item
При необходимости для автоматически ускоряемой ступени ДЗ можно ввести контроль от БНН программныv ключом \textbf{«БНН АУ ДЗ»}.

\item
Функциональная схема логики АУ ДЗ приведена на рисунке \ref{pict_dz:AU_DZ}.

\vspace{2mm}
\begin{figure}[H]
\centering
\includegraphics[width=1\textwidth,height=1\textheight,keepaspectratio]{\fbpath/03.01 ЛВ/561.0021 ДЗ/img/AU_DZ.pdf}
\captionof{figure}{Функциональная схема пуска АУ ДЗ}\label{pict_dz:AU_DZ}
\end{figure}

\item
Уставки автоматического ускорения ДЗ (АУ ДЗ) представлены в таблице \ref{app_set:dz}.


\end{enumerate}
%%%%%%%%%%%%%%%%%%%%%%%%%%%%%%%%%%%%%%%%%%%%%%%%%% КОНЕЦ АУ ДЗ %%%%%%%%%%%%%%%%%%%%%%%%%%%%%%%%%%%%%%%%%%%%%%%%%%


%%%%%%%%%%%%%%%%%%%%%%%%%%%%%%%%%%%%%%%%%%%%%%%%%% ОУ ДЗ %%%%%%%%%%%%%%%%%%%%%%%%%%%%%%%%%%%%%%%%%%%%%%%%%%
\item\label{sec:secou} Оперативное ускорение дистанционной защиты (ОУ ДЗ)
\begin{enumerate}

\item
В устройстве реализовано два логических узла оперативного ускорения ДЗ: ОУ1 ДЗ и ОУ2 ДЗ. Ввод в работу ОУ ДЗ осуществляется от внешнего сигнала \textit{«ОУ ДЗ: Ввод»}. Выбор оперативно ускоряемой ступени ДЗ выполняется с помощью программных переключателей \textbf{«Ступень ОУ ДЗ-ФФ»} и \textbf{«Ступень ОУ ДЗ-ФЗ»}. Время срабатывания ОУ ДЗ задается уставкой \textbf{«Тср ОУ ДЗ»}.

\item
Переключатель \textbf{«Направл. ОУ ДЗ»} задаёт направленность ОУ, что позволяет, например, организовать резервирование ДЗШ на «своей» подстанции. В зависимости от положения переключателя \textbf{«Направл. ОУ ДЗ»} действие ОУ ДЗ осуществляется от прямонаправленного либо обратнонаправленного РС выбранной ступени ДЗ-ФФ и ДЗ-ФЗ.

\item
Кроме того, для обеспечения надежного действия ОУ ДЗ, направленного назад, при КЗ на шинах, предусмотрен подхват сработанного состояния РС от ненаправленного РС. Подхват вводится программным ключами \textbf{«Подхват РС-ФФ ОУ ДЗ»} и \textbf{«Подхват РС-ФЗ ОУ ДЗ»}. Возврат схемы подхвата в исходное состояние происходит только после возврата ненаправленного РС.

\item
Предусмотрена возможность автоматического ввода ОУ ДЗ при приеме внешнего сигнала о выводе основной защиты. Автоматический ввод ОУ ДЗ разрешается программным ключом \textbf{«Ввод ОУ ДЗ при выводе осн. защиты»}. Введенное ускорение может быть оперативно выведено от внешнего сигнала \textit{«ОУ ДЗ: Вывод авт. введенного ОУ»}.

\item
При необходимости для оперативно ускоряемой ступени ДЗ можно ввести контроль от БК и БНН программными переключателями \textbf{«БК ОУ ДЗ»} и \textbf{«БНН ОУ ДЗ»}.

\item
Функциональная схема логики ОУ ДЗ приведена на рисунке \ref{pict_dz:OU_DZ}. Схема логики ОУ2 ДЗ выполнена аналогично ОУ1 ДЗ.

\item
Уставки оперативного ускорения ДЗ (ОУ ДЗ) представлены в таблице \ref{app_set:dz}.


\end{enumerate}
%%%%%%%%%%%%%%%%%%%%%%%%%%%%%%%%%%%%%%%%%%%%%%%%%% КОНЕЦ ОУ ДЗ %%%%%%%%%%%%%%%%%%%%%%%%%%%%%%%%%%%%%%%%%%%%%%%%%%


%%%%%%%%%%%%%%%%%%%%%%%%%%%%%%%%%%%%%%%%%%%%%%%%%% БСТО %%%%%%%%%%%%%%%%%%%%%%%%%%%%%%%%%%%%%%%%%%%%%%%%%%
\item\label{sec:secbsto} Блокировка при протекании сквозного тока (БСТО)
\begin{enumerate}

\item
Для предотвращения неправильной работы защиты предусмотрена блокировка при протекании через трансформаторы тока двух выключателей сквозного тока внешнего КЗ, вызывающего неодновременное насыщение трансформаторов. Блокировка вводится отдельными программными ключами для первых ступеней ДЗ-ФЗ-1 и ДЗ-ФФ-1 (\textbf{«БСТО ДЗ-ФФ-1»} и \textbf{«БСТО ДЗ-ФЗ-1»} соответственно), а также для обоих режимов оперативного ускорения (\textbf{«БСТО ОУ1 ДЗ»} и \textbf{«БСТО ОУ2 ДЗ»}).

\end{enumerate}
%%%%%%%%%%%%%%%%%%%%%%%%%%%%%%%%%%%%%%%%%%%%%%%%%% КОНЕЦ БСТО %%%%%%%%%%%%%%%%%%%%%%%%%%%%%%%%%%%%%%%%%%%%%%%%%%

\item
Дистанционная защита может быть оперативно выведена из работы от внешнего сигнала «Вывод ДЗ», назначенного на соответствующий дискретный вход или функциональную клавишу. Выводятся все ступени ДЗ, логика оперативного ускорения, оперативного ускорения без выдержки времени и автоматического ускорения.

\begin{figure}[H]
\centering
\includegraphics[width=0.95\textwidth,height=0.95\textheight,keepaspectratio]{\fbpath/03.01 ЛВ/561.0021 ДЗ/img/OU_DZ.pdf}
\captionof{figure}{Функциональная схема логики ОУ ДЗ}\label{pict_dz:OU_DZ}
\end{figure}

\end{enumerate}
